\documentclass[10pt,a4paper]{amsart}
\usepackage[margin=19mm]{geometry}
\usepackage{amsmath,amssymb,mathtools}
\usepackage[scr=rsfs,cal=euler]{mathalfa}
\usepackage{xfrac}
\usepackage[unicode,hidelinks,bookmarks=false]{hyperref}
\newcommand{\ie}{i.e.\ }
\newcommand{\cf}{cf.\ }
\newcommand{\resp}{respectively\ }
\newcommand{\Subsection}{\subsection}
\providecommand{\nobreakdash}{\nobreak\hbox{-}}
\setlength{\emergencystretch}{2em}
\multlinegap0pt
\usepackage{bm}
\usepackage{bbm}
\usepackage{dsfont}
\usepackage{xspace}
\usepackage{cancel}
\usepackage{mathtools}
\usepackage{amsthm}
\makeatletter
\@ifundefined{theorem}{\newtheorem{theorem}{Theorem}}{}
\@ifundefined{lemma}{\newtheorem{lemma}[theorem]{Lemma}}{}
\makeatother
\newcommand{\Id}{\operatorname{Id}}
\newcommand{\cB}{\mathcal B}
\newcommand{\cE}{\mathcal E}
\newcommand{\inner}[2]{\left\langle #1,#2\right\rangle}
\title[Pinching and Tensorial Rigidity for Ergodicity of Frame Flow]{Pinching and Tensorial Rigidity for Ergodicity of Frame Flows}
\author{Heng Zhang, Shuhao Zhang}
\date{Source: arXiv:2608.09600v1 (CC BY 4.0)}
\begin{document}
\maketitle

\noindent\emph{Source and licence.} Pinching and Tensorial Rigidity for Ergodicity of Frame Flows. Version arXiv:2608.09600v1 (CC BY 4.0), distributed under the Creative Commons Attribution 4.0 International licence, as recorded in the arXiv licence metadata for this exact version and shown on the abstract page https://arxiv.org/abs/2608.09600v1. Full rights notice: licenses/arxiv-2608.09600-CC-BY-4.0.txt.\par\smallskip

\noindent\textit{Editorial scope.} Counted unit: the Lemma whose source label is \texttt{lem:tangent-twist-residual} in the article (source0.tex lines 942--949, source label \texttt{lem:tangent-twist-residual}), delivered together with the article's own complete original proof of it (source0.tex lines 951--1128, one \texttt{proof} environment of 178 lines). No mathematics is written, repaired or re-derived: statement and proof are the article's own text line for line. Required context retained uncounted: lem:BBK (lines 844--855); eq:BK-simple-form (lines 868--871). Every definition, display and label of those lines is retained unchanged. Omitted from this package and not redistributed: 1--843, 856--867, 872--941, 950--950, 1129--4516. Nothing inside a retained excerpt is omitted or altered: every retained body is delivered exactly as the article's lines are written, so no omission receipt of any kind accompanies this package. Citations: the retained excerpts carry no citation command at all, so no bibliography entry of the article is transcribed and no reference is redistributed. Reference adaptation: none was needed and none was made; every reference inside the delivered unit resolves inside it, so no unresolved reference or citation is produced. Printed slips kept literal: no printed word, symbol, hypothesis or number of the counted statement or of its proof was altered. Layout and class: the article's own class is \texttt{amsart} with packages including amsmath, amssymb, amsfonts, amsthm, mathtools, mathrsfs, bm, enumitem, array, tabularx, microtype, hyperref, cleveref; the wrapper is the lane's pinned \texttt{amsart} editorial wrapper at 10pt on A4, which restates the theorem-like environments the retained text needs and adds the article's own macro definitions that the retained text needs (\texttt{\textbackslash Id}, \texttt{\textbackslash cB}, \texttt{\textbackslash cE}, \texttt{\textbackslash inner}), with extra wrapper packages (bm, bbm, dsfont, xspace, cancel, mathtools). The delivered numbering is the unit's own. Assets: no retained span contains a figure environment, a graphics-inclusion command, a PDF-page inclusion, a table or a TikZ picture; no image file is copied into this package, the package declares no asset dependency and no image file is redistributed with it. Licence: version arXiv:2608.09600v1 of this article is distributed under the Creative Commons Attribution 4.0 International licence, as recorded in the arXiv licence metadata for this exact version and shown on the abstract page https://arxiv.org/abs/2608.09600v1. A full rights notice is delivered at licenses/arxiv-2608.09600-CC-BY-4.0.txt. Characters: every retained excerpt is pure ASCII, so the delivered engine, XeLaTeX with its default Latin Modern text font, reports no missing character.
\begin{lemma}\label{lem:BBK}
Let $\cB$ be an algebraic curvature operator satisfying
$$
  a\leq\sec_{\cB}\leq b.
$$
Set $\widehat{\cB}:=\cB-\frac{a+b}{2}I$.  Then, for arbitrary vectors
$u,v,w,z$,
\begin{equation}\label{eq:BBK-vectors}
  |\widehat{\cB}(u,v,w,z)|
  \leq\frac23(b-a)|u|\,|v|\,|w|\,|z|.
\end{equation}
\end{lemma}

\begin{equation}\label{eq:BK-simple-form}
  |\cB(\xi,\eta)|
  \leq\frac23(b-a)|\xi|\,|\eta|.
\end{equation}

\begin{lemma}\label{lem:tangent-twist-residual}
If $0\leq\sec_{\cE}\leq1$, then
\begin{equation*}
  W_{\cE}(K)
  \geq-\frac14\bigl(N_{\perp}+(n-1)N_{\parallel}\bigr)
  -\frac23\sqrt{(n-2)N_{\perp}\,k(k+n-2)N}.
\end{equation*}
\end{lemma}

\begin{proof}
Fix $x\in S(V)$ and write
$$
  F=\phi x+F^\perp,
  \qquad
  F^\perp\perp x.
$$
Since the sum in the definition of $W_{\cE}(K)$ is independent of the
fixed ambient orthonormal basis, we may, when $F^\perp\neq0$, choose
$$
  e_0=x,
  \qquad
  e_1=\frac{F^\perp}{|F^\perp|},
  \qquad
  e_2,\ldots,e_{n-1}\perp\{x,F^\perp\}.
$$
These vectors are kept fixed as ambient vectors when differentiating the
component functions $F_r=\inner{F}{e_r}$. Put
$$
  z_r:=\nabla^S F_r(x)\in x^\perp,
  \qquad
  \beta_r=x\wedge z_r.
$$

Consider the Jacobi form on $x^\perp$,
$$
  \mathsf J_x^{\cE}(a,b)
  :=\cE(x\wedge a,x\wedge b).
$$
For $a\in x^\perp$ the sectional-curvature assumption gives
$$
  0\leq\mathsf J_x^{\cE}(a,a)\leq |a|^2,
$$
and hence
$$
  0\preceq\mathsf J_x^{\cE}\preceq\Id_{x^\perp}.
$$
Consequently, for all $a,z\in x^\perp$,
\begin{align}
  \mathsf J_x^{\cE}(z,z)-\mathsf J_x^{\cE}(a,z)
  &=
  \mathsf J_x^{\cE}(z-a/2,z-a/2)
  -\frac14\mathsf J_x^{\cE}(a,a) \geq-\frac14|a|^2.
  \label{eq:jacobi-complete-square}
\end{align}

We shall use the orthogonal-simple-form consequence
\eqref{eq:BK-simple-form} of Lemma~\ref{lem:BBK}, with $a=0$ and $b=1$.

We now estimate the summands
$$
  Q_r:=\cE(\beta_r,\beta_r)-\cE(\alpha_r,\beta_r).
$$
For $r=0$,
$$
  \alpha_0=F\wedge x=-x\wedge F^\perp,
$$
and therefore
$$
  Q_0
  =\mathsf J_x^{\cE}(z_0,z_0)
   -\mathsf J_x^{\cE}(-F^\perp,z_0)
  \geq-\frac14|F^\perp|^2
$$
by \eqref{eq:jacobi-complete-square}. Similarly,
$$
  \alpha_1=F\wedge e_1=\phi\,x\wedge e_1,
$$
because $F^\perp$ is parallel to $e_1$, and hence
$$
  Q_1
  =\mathsf J_x^{\cE}(z_1,z_1)
   -\mathsf J_x^{\cE}(\phi e_1,z_1)
  \geq-\frac14\phi^2.
$$

For $r\geq2$ we have
$$
  \alpha_r=\phi\,x\wedge e_r+F^\perp\wedge e_r.
$$
Thus
\begin{align*}
  Q_r
  ={}&
  \mathsf J_x^{\cE}(z_r,z_r)
  -\mathsf J_x^{\cE}(\phi e_r,z_r)
  -\cE(F^\perp\wedge e_r,\beta_r).
\end{align*}
The first two terms are bounded below by $-\phi^2/4$. Moreover,
$$
  \inner{F^\perp\wedge e_r}{x\wedge z_r}
  =
  \inner{F^\perp}{x}\inner{e_r}{z_r}
  -\inner{F^\perp}{z_r}\inner{e_r}{x}
  =0.
$$
Hence the two simple forms $F^\perp\wedge e_r$ and $\beta_r$ are
orthogonal, and \eqref{eq:BK-simple-form} gives
$$
  |\cE(F^\perp\wedge e_r,\beta_r)|
  \leq\frac23|F^\perp\wedge e_r|\,|\beta_r|
  =\frac23|F^\perp|\,|\beta_r|.
$$
It follows that
$$
  Q_r\geq-\frac14\phi^2
  -\frac23|F^\perp|\,|\beta_r|,
  \qquad r=2,\ldots,n-1.
$$

Summing the estimates for $r=0,\ldots,n-1$ yields the pointwise bound
\begin{equation}\label{eq:tangent-twist-pointwise}
  \sum_{r=0}^{n-1}
  \bigl\{\cE(\beta_r,\beta_r)-\cE(\alpha_r,\beta_r)\bigr\}
  \geq
  -\frac14\bigl(|F^\perp|^2+(n-1)\phi^2\bigr)
  -\frac23|F^\perp|
     \sum_{r=2}^{n-1}|\beta_r|.
\end{equation}
The same inequality holds when $F^\perp=0$ by continuity.

Multiplying by $c_{n,k}$ and integrating over $S(V)$, the first two
terms give
$$
  -\frac14\bigl(N_{\perp}+(n-1)N_{\parallel}\bigr).
$$
For the remaining term, Cauchy--Schwarz in the variables $(x,r)$ gives
\begin{align*}
  c_{n,k}\int_{S(V)}
  |F^\perp|\sum_{r=2}^{n-1}|\beta_r|\,d\sigma
  &\leq
  \left(
    c_{n,k}\int_{S(V)}
    \sum_{r=2}^{n-1}|F^\perp|^2\,d\sigma
  \right)^{1/2}\cdot
  \left(
    c_{n,k}\int_{S(V)}
    \sum_{r=2}^{n-1}|\beta_r|^2\,d\sigma
  \right)^{1/2}                                      \\
  &\leq
  \sqrt{(n-2)N_{\perp}}\,
  \left(
    c_{n,k}\int_{S(V)}
    \sum_{r=0}^{n-1}|\beta_r|^2\,d\sigma
  \right)^{1/2}.
\end{align*}
Although the adapted basis was chosen pointwise, the full sum
$\sum_r|\beta_r|^2$ is basis-independent. Using any fixed ambient
orthonormal basis,
$$
  \sum_{r=0}^{n-1}|\beta_r|^2
  =\sum_{r=0}^{n-1}|\nabla^S F_r|^2
  =|\nabla^S F|^2.
$$
Each $F_r$ is a spherical harmonic of degree $k$, since the corresponding
component of the homogeneous polynomial $F$ is harmonic. Therefore, with
$$
  \lambda_k=k(k+n-2),
$$
integration by parts on $S(V)$ gives
$$
  c_{n,k}\int_{S(V)}\sum_r|\beta_r|^2\,d\sigma
  =
  c_{n,k}\int_{S(V)}|\nabla^S F|^2\,d\sigma
  =
  \lambda_k c_{n,k}\int_{S(V)}|F|^2\,d\sigma
  =
  k(k+n-2)N.
$$
Substitution into \eqref{eq:tangent-twist-pointwise} proves
$$
  W_{\cE}(K)
  \geq-\frac14\bigl(N_{\perp}+(n-1)N_{\parallel}\bigr)
  -\frac23
   \sqrt{(n-2)N_{\perp}\,k(k+n-2)N},
$$
as claimed.
\end{proof}

\end{document}
